% s03-2-2-construction.tex —— 3.2.2 与 3.2.3 小节：AB 的构造与逐点等价（印刷页 86–91）

\subsection{AB 的构造}
\label{sec:ab-construction}

本小节说明如何在 $\FactModCat_{x_0}$ 中构造对的态射
\[
\mathrm{AB}\colon
(\Whit^{!}(G),\Whit^{!}(\Fl_G))\longrightarrow
(\Rep(\check G),\QCoh(\check{\mathfrak{n}}/\check B)).
\]

我们以对
\[
(\Whit^{!}(G),\Whit^{!}(LG)^{\infty}_{2})
\]
为中介。这里
\[
\Whit^{!}(LG)^{\infty}_{2}=\Whit^{!}(LG)^{LN L^+T},
\]
其中 $LN L^+T$ 指 $x_0$ 处的因子化 $L^+G$-模（ind-）群。

有函子
\[
(\Whit^{!}(G),\Whit^{!}(\Fl_G))\longrightarrow
(\Whit^{!}(G),\Whit^{!}(LG)^{\infty}_{2}),
\]
它是对逐点版本（该版本是等价）
\[
\Whit^{!}(\Fl_G)_{x_0}\xrightarrow{\ \sim\ }\Whit^{!}(LG)^{\infty}_{2,x_0}
\]
的直接升级；后者出现在 \cite[推论 6.2.3]{Ras16} 中（更精确地说，是作为该函子的逆，
即连续左伴随出现）。

我们将把 \cite[附录 B]{CCF+24} 的构造（该处把因子化范畴
$\Whit^{!}(LG)^{\infty}_{2}$ 与 $\IndCoh(\LS^{\mer}_{\check B}\times_{\LS^{\mer}_{\check T}}
\LS^{\reg}_{\check B})$ 联系起来）改造成因子化模范畴的情形，以构造映射
\[
(\Whit^{!}(G),\Whit^{!}(LG)^{\infty}_{2})\longrightarrow
(\Rep(\check G),\QCoh(\check{\mathfrak{n}}/\check B)).
\]
\srcpage{87}

在几何侧，考虑函子
\[
\Whit^{!}(LG)^{\infty}_{2}\longrightarrow
\Whit^{!}(T)\simeq D\text{-}\mathrm{mod}_{1/2}(\Gr_T),
\]
它由
\begin{align}
\Whit^{!}(LG)^{\infty}_{2}
&\subset D\text{-}\mathrm{mod}_{1/2}(LG)^{LN L^+T}\label{eq:3.2.1}\\
&\simeq \bigl(D\text{-}\mathrm{mod}_{1/2}(LG)\otimes
D\text{-}\mathrm{mod}_{1/2}(\Gr_T)\bigr)^{LB}\label{eq:3.2.2}\\
&\to D\text{-}\mathrm{mod}_{1/2}(LG)\otimes D\text{-}\mathrm{mod}_{1/2}(\Gr_T)
\xrightarrow{1^{!}_{LG}} D\text{-}\mathrm{mod}_{1/2}(\Gr_T)\label{eq:3.2.3}\\
&\xrightarrow{\ \text{移位}\ } D\text{-}\mathrm{mod}_{1/2}(\Gr^T).\label{eq:3.2.4}
\end{align}
给出。
它具有对的态射结构
\[
(\Whit^{!}(G),\Whit^{!}(LG)^{\infty}_{2})\longrightarrow
(\Whit^{!}(T),\Whit^{!}(T)),
\]
其中第一个分量 $\Whit^{!}(G)\to \Whit^{!}(T)$ 是 \cite{CCF+24} 中定义的 Jacquet 函子
$J^{\Whit^{!}}$。

该映射经
\[
(\Whit^{!}(G),\Whit^{!}(LG)^{\infty}_{2})\longrightarrow
\Omega\text{-}\mathrm{FactModCat}(\Whit^{!}(T),\Whit^{!}(T))
\]
分解。把 $\Omega\text{-}\mathrm{FactModCat}(\Whit^{!}(T))$ 沿
$\Whit^{!}(G)\to \Omega\text{-}\mathrm{FactModCat}(\Whit^{!}(T))$ 的因子化限制
（滥用记号地）记作 $(\Whit^{!}(G),\Omega\text{-}\FactMod(\Whit^{!}(T)))$，即得态射
\[
(\Whit^{!}(G),\Whit^{!}(LG)^{\infty}_{2})\longrightarrow
(\Whit^{!}(G),\Omega\text{-}\FactMod(\Whit^{!}(T))),
\]
\srcpage{88}
它在限制到 $x_0\in \Ran_{x_0}$ 时是等价。

在谱侧，可考虑函子
\[
\QCoh(\check{\mathfrak{n}}/\check B)\longrightarrow \Rep(\check T),
\]
它由沿
\[
\LS_{\check B}(D^\times)\times_{\LS_{\check T}(D^\times)}\LS_{\check T}(D)
\longrightarrow \LS_{\check T}(D)
\]
的推前给出。它具有对的态射结构
\[
(\Rep(\check G),\QCoh(\check{\mathfrak{n}}/\check B))\longrightarrow
(\Rep(\check T),\Rep(\check T)),
\]
其中第一个分量 $\Rep(\check G)\to \Rep(\check T)$ 由限制到 $\Rep(\check B)$ 再取
$\check N$-不变量给出。

该映射经
\[
(\Rep(\check G),\QCoh(\check{\mathfrak{n}}/\check B))\longrightarrow
\Omega^{\spec}\text{-}\FactMod(\Rep(\check T),\Rep(\check T))
\]
分解，后者又经
\[
(\Rep(\check G),\QCoh(\check{\mathfrak{n}}/\check B))\longrightarrow
(\Rep(\check G),\Omega^{\spec}\text{-}\FactMod(\Rep(\check T)))
\]
分解。该态射在限制到 $x_0\in \Ran_{x_0}$ 时是等价。此外，它在第二个分量的最终余连通
部分上也是等价，因为（$\FactAlgCat$ 中的）因子化函子
\[
\IndCoh\bigl(\LS^{\mer}_{\check B}\times_{\LS^{\mer}_{\check T}}
\LS^{\reg}_{\check B}\bigr)\longrightarrow
\Omega^{\spec}\text{-}\FactMod(\Rep(\check T))
\]
\srcpage{89}
由命题~\ref{prop:approx-ABC} 也是等价。

可以等同
\[
(\Whit^{!}(G),\Omega\text{-}\FactMod(\Whit^{!}(T)))
\simeq (\Rep(\check G),\Omega^{\spec}\text{-}\FactMod(\Rep(\check T))),
\]
因为 $\Omega$ 与 $\Omega^{\spec}$ 的匹配来自图
\[
\begin{tikzcd}[column sep=huge,row sep=large]
\Whit^{!}(G) \arrow[r,"\mathrm{CS}_G"] \arrow[d] & \Rep(\check G) \arrow[d]\\
\Whit^{!}(T) \arrow[r,"\mathrm{CS}_T"] & \Rep(\check T).
\end{tikzcd}
\]

最后，我们想构造所承诺的态射
\[
\mathrm{AB}\colon
(\Whit^{!}(G),\Whit^{!}(\Fl_G))\longrightarrow
(\Rep(\check G),\QCoh(\check{\mathfrak{n}}/\check B)),
\]
办法是说明 $\Whit^{!}(\Fl_G)$ 的紧生成元被送到
$\Omega^{\spec}\text{-}\FactMod(\Rep(\check T))$（等价地，$\Rep(\check T)$）中的
最终余连通对象。麻烦在于 $\Whit^{!}(\Fl_G)$ 大概不是紧生成的。然而
$\Whit^{!}(\Fl_G)$ 的可熔性救了场：因为
\[
\Whit^{!}(\Fl_G)_{\QLisse}\subset \Whit^{!}(\Fl_G)
\]
紧生成，且容易验证其紧生成元被送到最终余连通对象。于是得到所要的对的态射。
\srcpage{90}

\subsection{AB 是逐点等价}
\label{sec:ab-pointwise}

注意有链
\begin{align}
\Whit^{!}(\Fl_G)_{x_0}
&\xrightarrow{\ \sim\ }\Whit^{!}(LG)^{\infty}_{2,x_0}
\xrightarrow{\ \sim\ }\Omega\text{-}\FactMod(\Whit^{!}(T))_{x_0}
\label{eq:3.2.5}\\
&\simeq \Omega^{\spec}\text{-}\FactMod(\Rep(\check T))_{x_0}
\xleftarrow{\ \sim\ }\QCoh(\check{\mathfrak{n}}/\check B)_{x_0},
\label{eq:3.2.6}
\end{align}
故
\[
\mathrm{AB}_{x_0}\colon \Whit^{!}(\Fl_G)_{x_0}
\xrightarrow{\ \sim\ }\QCoh(\check{\mathfrak{n}}/\check B)_{x_0}
\]
是范畴的等价。由于 $\Whit^{!}(\Fl_G)$ 与 $\QCoh(\check{\mathfrak{n}}/\check B)$ 都可熔，
可以断定 AB 是 $\FactModCat_{x_0}$ 中的等价，这就证明了定理~\ref{thm:AB}。

不过，让我们再作一个平凡但后文有用的观察。在 \cite{AB09} 中，作者们（在重整化意义下）
构造了范畴的等价
\[
\mathrm{AB}^{\mathrm{classical}}\colon
\Whit^{!}(\Fl_G)_{x_0}\simeq \QCoh(\check{\mathfrak{n}}/\check B)_{x_0}.
\]
$\mathrm{AB}_{x_0}$ 与 $\mathrm{AB}^{\mathrm{classical}}$ 可能同构也可能不同构，但由于
两者都是等价，可以定义
\[
\alpha=\mathrm{AB}_{x_0}\circ \mathrm{AB}^{\mathrm{classical},-1}
\]
为 $\QCoh(\check{\mathfrak{n}}/\check B)_{x_0}$ 的自同构，
\[
\beta=\mathrm{AB}^{\mathrm{classical},-1}\circ \mathrm{AB}_{x_0}
\]
为 $\Whit^{!}(\Fl_G)_{x_0}$ 的自同构。于是
\srcpage{91}
\[
\mathrm{AB}_{x_0}=\alpha\circ \mathrm{AB}^{\mathrm{classical}},
\]
以及
\[
\mathrm{AB}_{x_0}=\mathrm{AB}^{\mathrm{classical}}\circ \beta .
\]
